\paragraph{Step 25 co-sourcing criterion.}
Let \(L\) carry a finite complex field \(\psi\) on the toy lattice.
The QM descent reads
\[
  \operatorname{Born}[\psi]=(\rho,j),
  \qquad
  \rho=|\psi|^2,\quad
  j=\operatorname{Im}(\overline\psi\,\nabla\psi),
\]
and the GR descent reads
\[
  T[\psi]=(T_{00},T_{0i}),
  \qquad
  T_{00}=|\nabla\psi|^2+V|\psi|^2,\quad
  T_{0i}=\operatorname{Re}(\overline{\nabla\psi}\,\Delta_+\psi).
\]
The co-sourcing unification criterion is that the same field \(\psi\)
on \(L\) sources both descent audits:
\[
  \operatorname{QM\ descent}=\operatorname{Born}[\psi],
  \qquad
  \operatorname{GR\ descent}=T[\psi].
\]

\paragraph{Computed result.}
For the co-sourced \(L\), both residuals are zero across all sampled
fields.  For the non-co-sourced control, the QM descent is still
Born\([\psi]\), but the GR descent is \(T[\psi']\) with
\(\psi'\ne\psi\); the maximum stress-energy residual is
\(0.3806026603775931\).  Thus \(L\) satisfies the co-sourcing
criterion and the control fails.
